\chapter{\B{Generative models at scale}{Modelos generativos a gran escala}}
\label{ch:generative}

\B{You run the service desk at a water and power company in Curaçao. By midday you have three jobs. The first is to sort a pile of messages into water, power and billing. The second is to draft a reply. The third is to translate a notice so customers can read it in another language.

You can pick the first tool your desk will test. Your team has time for one small trial. The duty officer approves changes to the official notice. Start with the job that takes most of your time and decide what a good result would look like. Then match the job to what each type of model can read and produce.}
{Diriges la mesa de atención de una empresa de agua y electricidad en Curazao. Al mediodía tienes tres tareas. Primero, ordenar mensajes de agua, luz y facturación. Después, preparar una respuesta. También debes traducir un aviso para que los clientes lo lean en otro idioma.

Tú eliges qué herramienta probará primero la mesa. Tu equipo tiene tiempo para una prueba pequeña. La persona de turno autoriza los cambios al aviso oficial. Empieza por la tarea que más tiempo te quita y define un buen resultado. Luego relaciona la tarea con lo que cada tipo de modelo puede leer y producir.}

\section{\B{Choose an objective and an architecture}{Elige un objetivo y una arquitectura}}

\B{For a draft reply, think of the model writing “We can help with your” and then choosing what to add. The words it has so far are the \term{prefix}. A \term{decoder-only} model uses a mask so it can look back at that prefix when it picks the next piece. This rule is called \term{causal attention}.

During practice, the saved text already gives the model each prefix. It can score guesses for many spots at the same time. The mask keeps the next piece hidden at each spot. When the model writes a new reply, it has to pick one piece before it knows the next prefix \citep{xiao,gpt3}.}
{Para preparar una respuesta, piensa en un modelo que escribe «Podemos ayudarte con tu» y luego elige qué añadir. Las palabras que lleva se llaman \term{prefijo}. Un modelo de \term{solo decodificador} usa una máscara para mirar el prefijo al elegir la siguiente pieza. La regla se llama \term{atención causal}.

Durante la práctica, el texto guardado ya da cada prefijo. El modelo puede puntuar sus intentos en muchos lugares a la vez. La máscara oculta la pieza siguiente en cada lugar. Cuando escribe una respuesta nueva, debe elegir una pieza antes de conocer el siguiente prefijo \citep{xiao,gpt3}.}

\B{For sorting, the model can read the whole message before it labels it. A \term{model with only an encoder} can use words on both sides of any piece. Give it “The water came back, but my bill is wrong.” It can use “bill is wrong” to sort the message into billing.

One way to train the encoder is to hide or change some pieces and ask it to guess what was there. BERT learned that task and also practised telling whether one sentence followed another in the training text. A part added to the encoder can learn to give each message one label. It is called a \term{classification head} \citep{bert}.}
{Para ordenar mensajes, el modelo puede leer el mensaje completo antes de etiquetarlo. Un \term{modelo de solo codificador} puede usar las palabras a ambos lados de una pieza. Dale «Volvió el agua, pero mi factura está mal». Puede usar «mi factura está mal» para poner el mensaje en facturación.

Una forma de entrenar el codificador es ocultar o cambiar algunas piezas y pedirle que adivine qué había. BERT aprendió esa tarea y también practicó decidir si una oración seguía a otra en el texto de entrenamiento. Una parte que se añade al codificador puede aprender a dar una etiqueta a cada mensaje. Se llama \term{cabeza de clasificación} \citep{bert}.}

\B{To translate a notice, you may want the tool to read it from start to finish before it writes the first word. In an \term{encoder--decoder} model, the encoder reads the notice. The decoder can look at the encoder's notes as it writes each piece of the translation. This link is called \term{cross-attention}. The decoder also looks at the translation it has written so far.

T5 practised filling gaps in text. For a similar exercise, replace “tomorrow” and “Orange Walk” with markers in a produce notice from Belize. The model reads the notice with its gaps and writes the missing parts \citep{t5}.}
{Para traducir un aviso, quizá quieras que la herramienta lo lea de principio a fin antes de escribir la primera palabra. En un modelo \term{codificador--decodificador}, el codificador lee el aviso. El decodificador puede mirar el registro que hizo el codificador mientras escribe cada pieza de la traducción. Este enlace se llama \term{atención cruzada}. El decodificador también mira lo que ya lleva escrito.

T5 practicó llenar huecos en el texto. Para un ejercicio parecido, cambia «mañana» y «Orange Walk» por marcas en un aviso de productos agrícolas de Belice. El modelo lee el aviso con sus huecos y escribe las partes que faltan \citep{t5}.}

\begin{figure}[htbp]\centering
\begin{tikzpicture}[node distance=5mm]
\foreach \y/\src/\body/\outcome in {
 0/{\B{Visible context}{Contexto visible}}/{\B{Bidirectional encoder}{Codificador bidireccional}}/{\B{Representations}{Representaciones}},
 -1.22/{\B{Output prefix}{Prefijo de salida}}/{\B{Causal decoder}{Decodificador causal}}/{\B{Next-token scores}{Puntuaciones del siguiente token}},
 -2.44/{\B{Source and prefix}{Origen y prefijo}}/{\B{Encode and cross-attend}{Codifica y aplica atención cruzada}}/{\B{Target scores}{Puntuaciones del objetivo}}}{
 \node[nodebox,fill=softgrey] (s) at (0,\y){\src};
 \node[nodebox] (m) at (4.55,\y){\body};
 \node[nodebox,fill=softgreen] (o) at (9.10,\y){\outcome};
 \draw[mainarr](s)--(m);\draw[mainarr](m)--(o);
}
\end{tikzpicture}
\caption{\B{The rows contrast which tokens are visible. An encoder can also feed a prediction head, and a decoder can classify by generating a label.}{Las filas contrastan qué tokens están visibles. Un codificador también puede alimentar una cabeza de predicción y un decodificador puede clasificar generando una etiqueta.}}
\label{fig:architectures}
\end{figure}

\section{\B{Multiply conditional probabilities}{Multiplica probabilidades condicionales}}

\B{A reply comes one piece at a time. The chance of “today” can change once the model has picked “ready”. A chance that depends on what has already happened is called a \term{conditional probability}. To find the chance of a whole reply, multiply the chance of each next piece along that path.

Write \(x\) for the prompt and \(y_1,\ldots,y_m\) for the pieces in the reply. At step \(t\), \(y_{<t}\) means the part already written. The vertical line in \(p_\theta(y_t\mid x,y_{<t})\) means “given”. The big \(\prod\) tells you to multiply all the chances.}
{Una respuesta sale pieza por pieza. La probabilidad de «hoy» puede cambiar cuando el modelo ya eligió «listo». Una probabilidad que depende de lo que ya pasó se llama \term{probabilidad condicional}. Para saber la probabilidad de toda una respuesta, multiplica la probabilidad de cada pieza siguiente en ese camino.

Escribe \(x\) para la instrucción y \(y_1,\ldots,y_m\) para las piezas de la respuesta. En el paso \(t\), \(y_{<t}\) es lo que ya se escribió. La barra en \(p_\theta(y_t\mid x,y_{<t})\) significa «dado». El símbolo grande \(\prod\) indica que hay que multiplicar las probabilidades.}
\begin{equation}
 p_\theta(y_1,\ldots,y_m\mid x)=
 \prod_{t=1}^{m}p_\theta(y_t\mid x,y_{<t}).
 \label{eq:chain}
\end{equation}
\B{Multiply the probability of each specified token given the preceding tokens. Include an end marker to score the probability of ending at that point. Changing the first token can change every distribution after it. A likelihood belongs to a particular model, tokenisation and context. The factual accuracy of the resulting text requires a separate check.}
{Multiplica la probabilidad de cada token dado lo anterior. Incluye una marca de fin para puntuar la probabilidad de terminar en ese punto. Cambiar el primer token puede alterar todas las distribuciones posteriores. Una verosimilitud corresponde a un modelo, tokenización y contexto específicos. La exactitud factual del texto requiere otra comprobación.}

\lab{\B{A short reply in Saint Vincent}{Una respuesta breve en San Vicente}}
\B{An assistant drafts a reply to a practice question about a produce shipment. Suppose the model gives \emph{ready}, \emph{today} and the end marker conditional probabilities \(0.6,0.5,0.8\) in a toy tokenizer. The complete reply has probability \(0.24\). The two-token prefix has probability \(0.30\). The gap comes from the chance of continuing after \emph{today}. The shipment clerk checks the actual shipment record. Hold the reply until that person confirms it.}
{Un asistente redacta una respuesta a una pregunta de práctica sobre un envío agrícola. Supón que el modelo da probabilidades condicionales \(0.6,0.5,0.8\) a \emph{listo}, \emph{hoy} y la marca de fin con un tokenizador didáctico. La respuesta completa tiene probabilidad \(0.24\). El prefijo de dos tokens tiene probabilidad \(0.30\). La diferencia proviene de la posibilidad de continuar después de \emph{hoy}. La persona a cargo de los envíos revisa el registro. Espera su confirmación antes de enviar la respuesta.}

\B{Multiply the first two chances to get the chance of the two pieces. Then multiply by the chance of ending to get the chance of the full reply.}
{Multiplica las dos primeras probabilidades para obtener la probabilidad de las dos piezas. Luego multiplica por la probabilidad de terminar para obtener la de la respuesta completa.}
\[0.6\times0.5=0.30,\qquad 0.30\times0.8=0.24.\]

\B{For a long reply, multiplying many small chances can give a number too tiny for a computer to store. You can score the reply another way. Press the \(\ln\) key on each probability, then add the results. This is the sum of their \term{logarithms}. It ranks replies in the same order as multiplying their chances. A longer reply usually has more terms to add. If you use a rule that adjusts the score for length, write it down before you compare replies.}
{En una respuesta larga, al multiplicar muchas probabilidades pequeñas puedes obtener un número demasiado pequeño para que una computadora lo guarde. Puedes puntuar la respuesta de otra forma. Pulsa \(\ln\) para cada probabilidad y suma los resultados. Esta es la suma de sus \term{logaritmos}. Ordena las respuestas igual que al multiplicar sus probabilidades. Una respuesta larga suele tener más términos que sumar. Si usas una regla que ajusta la puntuación según la longitud, anótala antes de comparar respuestas.}

\section{\B{Allocate data and compute}{Distribuye datos y cómputo}}

\B{Your service desk can try a tool that already exists. Building a model from the start takes much more work. Every piece of training text has to pass through the model's settings. More settings or more pieces mean more steps for the computer.

\moremath
Call the number of settings \(N\) and the number of training pieces \(D\). A computer adds and multiplies numbers to train the model. We count those steps as operations and call the total \(C\). The estimate below uses about six steps per setting and piece for a dense Transformer. In a dense model, the main groups of weights are used for every piece \citep{hoffmann}.}
{Tu mesa de atención puede probar una herramienta que ya existe. Crear un modelo desde cero lleva mucho más trabajo. Cada pieza del texto de práctica pasa por los ajustes del modelo. Más ajustes o más piezas significan más pasos para la computadora.

\moremath
Llama \(N\) al número de ajustes y \(D\) al número de piezas de entrenamiento. Una computadora suma y multiplica para entrenar el modelo. Contamos esos pasos como operaciones y llamamos \(C\) al total. La estimación usa unos seis pasos por ajuste y pieza para un Transformer denso. En un modelo denso se usan los grupos principales de pesos para cada pieza \citep{hoffmann}.}
\begin{equation}
 C\approx 6ND .
 \label{eq:compute}
\end{equation}
\B{Multiply parameters by processed tokens and by about six operations per parameter-token pair. For \(N=10^8\) and \(D=2\times10^9\), the approximation gives \(1.2\times10^{18}\) operations. \(2\times10^8\) parameters and \(10^9\) tokens give the same estimate. Those plans can give different results. Long attention windows, embeddings, padding and device efficiency can move actual costs away from this estimate. Wall-clock time needs a throughput measurement.}
{Multiplica los parámetros por los tokens procesados y por unas seis operaciones por par parámetro-token. Para \(N=10^8\) y \(D=2\times10^9\), la aproximación da \(1.2\times10^{18}\) operaciones. Con \(2\times10^8\) parámetros y \(10^9\) tokens obtienes la misma estimación. Los planes pueden dar resultados distintos. Las ventanas de atención largas, las representaciones, el relleno y la eficiencia del equipo pueden alejar los costos reales de la estimación. Para estimar el tiempo transcurrido hace falta medir el rendimiento.}

\B{Researchers sometimes fit a curve to see how the loss changes as the model or its training set grows. The curve is a way to describe what happened in those tests. The equation gives one form for that curve.}
{A veces se ajusta una curva para ver cómo cambia la pérdida cuando crece el modelo o su conjunto de entrenamiento. La curva describe lo que pasó en esas pruebas. La ecuación da una forma posible de la curva.}
\begin{equation}
 L(N,D)\approx L_\infty + A N^{-\alpha}+B D^{-\beta}.
 \label{eq:scaling}
\end{equation}
\B{Here \(L\) is the loss, \(L_\infty\) a fitted floor, \(A,B\) positive fitted amplitudes and \(\alpha,\beta\) positive exponents. Increasing either \(N\) or \(D\) shrinks the corresponding term, with diminishing changes along the curve. If \(\alpha=1/2\), quadrupling \(N\) halves the model-size term while leaving the data term fixed. A new language mixture or repeated passes over a small corpus can break the fit's assumptions. Count unique material as well as processed tokens \citep{hoffmann}.}
{Aquí \(L\) es la pérdida, \(L_\infty\) un piso ajustado, \(A,B\) amplitudes positivas ajustadas y \(\alpha,\beta\) exponentes positivos. Aumentar \(N\) o \(D\) reduce el término correspondiente, con cambios menores a lo largo de la curva. Si \(\alpha=1/2\), cuadruplicar \(N\) reduce a la mitad el término del tamaño y mantiene fijo el de datos. Una mezcla nueva de idiomas o repeticiones sobre un corpus pequeño pueden romper los supuestos del ajuste. Cuenta material único además de tokens procesados \citep{hoffmann}.}

\B{Suppose a desk in Suriname has trouble with a reply about a trade form. Check the source first. Is the form missing? Was the scan hard to read? Did the tool pick the wrong file or mistranslate a line? Try the same question with a clear copy before paying for a bigger model. Use some forms the model has never seen and check each language separately.}
{Supón que una mesa de atención en Surinam tiene problemas con una respuesta sobre un formulario de comercio. Revisa primero la fuente. ¿Falta el formulario? ¿Cuesta leer el escaneo? ¿La herramienta eligió otro archivo o tradujo mal una línea? Prueba la misma pregunta con una copia clara antes de pagar por un modelo más grande. Usa algunos formularios que el modelo no haya visto y comprueba cada idioma por separado.}

\section{\B{Split the work across devices}{Reparte el trabajo entre dispositivos}}

\B{If your team has a pile of messages to sort, each person can work on a different part of the pile. A model can also share the work. In \term{data parallelism}, several devices hold copies of the same model. Each copy practises on a different set of examples. The devices combine what they learned about how to change the settings, then take one shared step.

\moremath
Let \(g\) count the copies. Each copy sees \(b\) text sequences in a small batch, called a \term{microbatch}. It saves the changes for \(a\) small batches before the shared step. \(B_{\mathrm{seq}}\) counts the sequences in all those batches.}
{Si tu equipo tiene una pila de mensajes para ordenar, cada persona puede trabajar con una parte. Un modelo también puede repartir el trabajo. En el \term{paralelismo de datos}, varios dispositivos guardan copias del mismo modelo. Cada copia practica con ejemplos distintos. Los equipos juntan lo que aprendieron sobre cómo cambiar los ajustes y dan un paso en común.

\moremath
Sea \(g\) el número de copias. Cada copia ve \(b\) secuencias de texto en un lote pequeño, llamado \term{microlote}. Guarda los cambios de \(a\) lotes pequeños antes del paso en común. \(B_{\mathrm{seq}}\) cuenta las secuencias de todos esos lotes.}
\begin{equation}
 B_{\mathrm{seq}}=g\,b\,a .
 \label{eq:batch}
\end{equation}
\B{Multiply the three counts. With \(4\) replicas, \(2\) sequences per replica and \(8\) accumulation steps, one update covers 64 sequences. If each contributes 512 unmasked loss positions, that is 32,768 scored tokens. Variable lengths or padding change the token count. For a per-token objective, add token penalties and divide by the total unmasked token count across replicas and accumulated microbatches.}
{Multiplica los tres conteos. Con \(4\) réplicas, \(2\) secuencias por réplica y \(8\) pasos de acumulación, una actualización abarca 64 secuencias. Si cada una aporta 512 posiciones de pérdida sin máscara, son 32\,768 tokens puntuados. Las longitudes variables o el relleno cambian el conteo. Para un objetivo por token, suma penalizaciones y divide por el total de tokens sin máscara entre réplicas y microlotes acumulados.}

\B{You can also share one large calculation. Each device handles part of it and then sends its result. That is \term{tensor parallelism}. Or you can put early model blocks on one device and later blocks on another. The first sends its results to the second, like desks that pass a form along. That is \term{pipeline parallelism}. You can also spread the stored settings across devices to save room. Sending results takes time, so try a small batch on the network you plan to use \citep{xiao}.}
{También puedes repartir un solo cálculo grande. Cada equipo hace su parte y luego envía el resultado. Eso es \term{paralelismo tensorial}. Otra opción es poner los primeros bloques del modelo en un equipo y los últimos en otro. El primero pasa sus resultados al segundo, como ventanillas que se pasan un formulario. Eso es \term{paralelismo de tubería}. También puedes repartir los ajustes guardados para ahorrar espacio. Enviar resultados toma tiempo. Prueba un lote pequeño en la red que piensas usar \citep{xiao}.}

\begin{figure}[htbp]\centering
\flow{\B{Replica 1\\Microbatch A}{Réplica 1\\Microlote A}}{\B{Gradients from A}{Gradientes de A}}{\B{Combine with gradients from B}{Combina con gradientes de B}}{\B{Apply a shared update}{Aplica un paso común}}
\caption{\B{A second replica processes B. The merge indicates the communication that keeps the replicas in sync.}{Una segunda réplica procesa B. La unión indica la comunicación que mantiene sincronizadas las réplicas.}}
\label{fig:parallel}
\end{figure}

\section{\B{Read longer text within a memory budget}{Lee textos largos con un presupuesto de memoria}}

\B{Put the same notice at the top and side of a square grid, one token per line. A cell shows where the model can compare two tokens. A longer notice makes the grid wider and taller. For \(n\) tokens, a full grid has \(n^2\) cells per head. A causal mask keeps \(n(n+1)/2\) of the pairs. Doubling the number of tokens gives four times as many cells in the full grid.

\term{FlashAttention} works on blocks of the grid and avoids storing it all in fast memory. Full attention still has to work with pairs of tokens, so that work grows roughly with the square of the text length \citep{flash}.}
{Pon el mismo aviso arriba y al lado de una cuadrícula, con un token por línea. Cada casilla indica dónde el modelo puede comparar dos tokens. Un aviso más largo hace crecer la cuadrícula en las dos direcciones. Con \(n\) tokens, la cuadrícula completa tiene \(n^2\) casillas por cabeza. Una máscara causal permite \(n(n+1)/2\) pares. Al duplicar el número de tokens, hay cuatro veces más casillas en la cuadrícula completa.

\term{FlashAttention} trabaja con bloques de la cuadrícula y evita guardarla entera en memoria rápida. La atención completa aún trabaja con pares de tokens, así que ese trabajo crece más o menos con el cuadrado del largo del texto \citep{flash}.}

\B{A sliding window restricts the prior positions each query can use in one layer. A 256-position window at each of 2,048 positions permits at most \(2{,}048\times256=524{,}288\) pairs if the current position counts. Early rows permit fewer pairs. Stacked layers can relay information over several windows. Some systems combine local and global attention. Others retrieve selected text before generation.}
{Una ventana deslizante restringe las posiciones anteriores que puede usar una consulta en una capa. Una ventana de 256 posiciones para cada una de 2\,048 permite como máximo \(2{,}048\times256=524{,}288\) pares si incluye la posición actual. Las primeras filas permiten menos pares. Capas sucesivas pueden transmitir información a través de varias ventanas. Algunos sistemas combinan atención local y global. Otros recuperan texto seleccionado antes de generar.}

\B{An extended positional scheme can let a model accept longer inputs. Check what it uses at that length. Move a dated ship notice between the beginning, middle and end of a long prompt, holding other content fixed. In a separate test, change one date in a distractor without moving the applicable notice. Measure how often retrieval finds the right record and how often the final answer uses it. These are separate outcomes \citep{lostmiddle}.}
{Un esquema posicional extendido puede permitir entradas más largas. Comprueba qué usa el modelo a esa longitud. Mueve un aviso marítimo fechado entre el principio, el medio y el final de una instrucción larga y mantén fijo el resto. En otra prueba, cambia una fecha en un distractor sin mover el aviso aplicable. Mide con qué frecuencia la recuperación encuentra el registro adecuado y con qué frecuencia la respuesta final lo usa. Son resultados distintos \citep{lostmiddle}.}

\lesson{\B{Choose your first trial}
{Elige tu primera prueba}}
\B{If sorting is the slow job, test whether the tool puts messages in the right pile. If replies are piling up, test a drafting model and have a teammate check the facts before sending. If translation is the bigger need, give the tool the whole notice. Ask a person who reads both languages to check the dates, place names and times. You decide whether the test earns a wider trial. If it gets the labels, facts or wording wrong, keep the current process while your team fixes and tests it again. The duty officer still signs off on any notice that goes out.}
{Si ordenar mensajes es lo que más demora, prueba si la herramienta los pone en la categoría correcta. Si se acumulan las respuestas, prueba un modelo que prepare borradores y pide que alguien revise los datos antes de enviarlos. Si la traducción hace más falta, dale el aviso completo. Pide a alguien que lea los dos idiomas que revise fechas, lugares y horas. Tú decides si vale la pena ampliar la prueba. Si falla en las etiquetas, los datos o la redacción, conserva el proceso actual mientras tu equipo lo corrige y vuelve a probar. La persona de turno sigue aprobando cada aviso que se publica.}

\checklist{
\item \B{Choose an objective and identify its visible context.}{Elegir un objetivo e identificar su contexto visible.}
\item \B{Separate an operations estimate, device memory and elapsed time.}{Distinguir una estimación de operaciones, memoria del equipo y tiempo transcurrido.}
\item \B{Design a test for use of evidence in a long prompt.}{Diseñar una prueba del uso de evidencia en una entrada larga.}}
\exercises{
\item \B{A reply uses three tokens with conditional probabilities \(0.4,0.5,0.5\), with the last token marking the end. Compute the complete and two-token prefix probabilities.}{Una respuesta usa tres tokens de probabilidades condicionales \(0.4,0.5,0.5\), con el último como marca de fin. Calcula las probabilidades de la respuesta completa y del prefijo de dos tokens.}
\item \B{Under Equation~\eqref{eq:compute}, hold \(C\) fixed and triple \(N\). How must \(D\) change? Give a reason measured time could still differ.}{En la ecuación~\eqref{eq:compute}, fija \(C\) y triplica \(N\). ¿Cómo debe cambiar \(D\)? Da una razón por la que el tiempo medido aún podría diferir.}
\item \B{Use three replicas, two sequences per replica and four accumulation steps. If all sequences score 256 positions, how many scored tokens does an update cover?}{Usa tres réplicas, dos secuencias por réplica y cuatro pasos de acumulación. Si cada secuencia puntúa 256 posiciones, ¿cuántos tokens puntuados abarca una actualización?}
\item \B{Why could a notice containing the answer still yield a wrong reply in a long prompt? Name a controlled test.}{¿Por qué un aviso con la respuesta podría generar una contestación errónea en una entrada larga? Indica una prueba controlada.}}
